\chapter{Sequence Diagram}
\label{ch:sequence}

\section{Purpose}

The sequence diagram models interaction between the participants over time,
showing the order in which messages pass between the engineer, the tool and
GENESYS for both translation directions.

\umlfigure{sequence-diagram}{Sequence diagram for both translation directions}{fig:sequence}

\section{Participants}

\begin{longtable}{L{0.20\textwidth}L{0.72\textwidth}}
	\caption{Sequence participants}\label{tab:sequence-participants}\\
	\toprule
	\textbf{Lifeline} & \textbf{Role} \\
	\midrule
	\endfirsthead
	\toprule
	\textbf{Lifeline} & \textbf{Role} \\
	\midrule
	\endhead
	User & The engineer. Initiates each translation and receives its result. \\
	\addlinespace
	Pentateuch & The translation tool. Performs extraction, abstraction, mapping
	and generation. Drawn as one lifeline because the internal division of
	responsibility is the subject of Figure~\ref{fig:block} rather than of this
	diagram. \\
	\addlinespace
	Genesys & The host application and its repository. Receives and supplies
	models. \\
	\bottomrule
\end{longtable}

\section{Message Sequence}

\begin{longtable}{L{0.06\textwidth}L{0.24\textwidth}L{0.24\textwidth}L{0.34\textwidth}}
	\caption{Messages}\label{tab:sequence-messages}\\
	\toprule
	\textbf{\#} & \textbf{From} & \textbf{To} & \textbf{Message} \\
	\midrule
	\endfirsthead
	\toprule
	\textbf{\#} & \textbf{From} & \textbf{To} & \textbf{Message} \\
	\midrule
	\endhead
	1 & User       & Pentateuch & Upload Code \\
	2 & Pentateuch & Genesys    & Import Model \\
	3 & Genesys    & Pentateuch & Export Model \\
	4 & Pentateuch & User       & Export Code \\
	\bottomrule
\end{longtable}

Messages 1 and 2 form the code-to-model direction: the engineer supplies source,
and the tool places the resulting model into GENESYS. Messages 3 and 4 form the
model-to-code direction: GENESYS supplies the model, and the tool returns
generated source.

\section{Design Decisions}

\paragraph{One lifeline for the tool.}
The tool is represented by a single lifeline rather than by its internal
components. This keeps the diagram at the level of interaction between the
engineer and the two systems, and leaves internal structure to
Figure~\ref{fig:block}.
